\documentclass{article}

\usepackage{amsmath}
\usepackage{amssymb}
\usepackage{float}
\usepackage{tikz}
\usepackage{xcolor}


\usepackage[utf8]{inputenc}
\usepackage[T1]{fontenc}

% \usepackage[polish]{babel}

\usepackage{listings}
\lstset{
    language=SQL,
    inputencoding=utf8x,
    extendedchars=\true,
    literate={ą}{{\k{a}}}1
             {Ą}{{\k{A}}}1
             {ę}{{\k{e}}}1
             {Ę}{{\k{E}}}1
             {ó}{{\'o}}1
             {Ó}{{\'O}}1
             {ś}{{\'s}}1
             {Ś}{{\'S}}1
             {ł}{{\l{}}}1
             {Ł}{{\L{}}}1
             {ż}{{\.z}}1
             {Ż}{{\.Z}}1
             {ź}{{\'z}}1
             {Ź}{{\'Z}}1
             {ć}{{\'c}}1
             {Ć}{{\'C}}1
             {ń}{{\'n}}1
             {Ń}{{\'N}}1
}

\newcommand\hlight[1]{\tikz[overlay, remember picture,baseline=-\the\dimexpr\fontdimen22\textfont2\relax]\node[rectangle,fill=blue!50,rounded corners,fill opacity = 0.2,draw,thick,text opacity =1] {$#1$};}

\title{Rozwiązywanie Układów Równań}
\date{2026}

\begin{document}
\maketitle

%%%%%%%%%%%%%%%%%%%%%%%%%%%%%%%%%%%%%%%%%%%%%%%%%%%%%%%%%%%%%%%%
\section{Rząd macierzy}
%===============================================================
\subsection{Definicja}
Maksymalny stopień wyznacznika z macierzy $A = [a_{ij}]_{n \times m}$,
różnego od zera nazywamy rzędem macierzy i oznaczamy $rk(A)$.

Dodatkowo przyjmujemy, że rząd macierzy zerowej jest równy zero.

%===============================================================
\subsection{Własności}
Rząd macierzy $A = [a_{ij}]_{n \times m}$ nie ulega zmianie w wyniku
następujących przekształceń:
\begin{itemize}
    \item transpozycji $rk(A) = rk(A^{T})$,
    \item pomnożenie dowolnego wiersza/kolumny macierzy przez 
        liczbę różną od zera,
    \item przestawienia wierszy/kolumn,
    \item dodania do wiersza/kolumny przemnożonego innego wiersza/kolumny,
    \item skreślenia lub dodania wiersza/kolumny skladajacego się 
        z samych zer.
\end{itemize}

%%%%%%%%%%%%%%%%%%%%%%%%%%%%%%%%%%%%%%%%%%%%%%%%%%%%%%%%%%%%%%%%
\section{Równanie liniowe}
%===============================================================
\subsection{Definicja}
Ogólna postać układu równań liniowych:
$$
    \left\{
    \begin{array}{l}
        a_{11}x_{1} + a_{12}x_{2} + \dots + a_{1n}x_{n} = b_{1}\\
        a_{21}x_{1} + a_{22}x_{2} + \dots + a_{2n}x_{n} = b_{2}\\
        \vdots\\
        a_{m1}x_{1} + a_{m2}x_{2} + \dots + a_{mn}x_{n} = b_{m}\\
    \end{array}
    \right.
$$
Który można zapisać w zwartej postaci jako:
$$
    Ax = b
$$
gdzie:
\begin{itemize}
    \item 
    $
        \left[
        \begin{array}{cccc}
            a_{11}  &   a_{12}  &   \cdots  &   a_{1n}\\
            a_{21}  &   a_{22}  &   \cdots  &   a_{2n}\\
            \vdots  &   \vdots  &   \ddots  &   \vdots\\
            a_{m1}  &   a_{m2}  &   \cdots  &   a_{mn}\\
        \end{array}
        \right]
    $ - macierz współczynników,
    \item
    $
    \left[
        \begin{array}{c}
            b_{1}\\
            b_{2}\\
            \vdots\\
            b_{m}\\
        \end{array}
        \right]
    $ - wektor wyrazów wolnych,
    \item
    $
        \left[
        \begin{array}{c}
            x_{1}\\
            x_{2}\\
            \vdots\\
            x_{n}\\
        \end{array}
        \right]
    $ - wektor niewiadomych,
\end{itemize}

Przypadek gdy wektor wyrazów wolnych utworzony jest z samych zer,
tzn. $Ax = \mathbf{0}$, nazywamy układem jednorodnym.

%%%%%%%%%%%%%%%%%%%%%%%%%%%%%%%%%%%%%%%%%%%%%%%%%%%%%%%%%%%%%%%%
\section{Twierdzenie Cramera}
%===============================================================
Jeżeli w układzie $Ax = b$, gdzie $A \in \mathbf{R}^{n \times n}$,
$x \in \mathbf{R}^{n}$ oraz $b \in \mathbf{R}^{n}$, a macierz 
współczynników ma niezerowy wyznacznik $|A| \neq 0$, to układ ma 
dokładnie jedno rozwiązanie, które dane jest wzorem Cramera:
$$
    x_{i} = \frac{W_{i}}{|A|}
$$
dla $i = 1, 2, \dots, n$.

Gdzie $W_{i}$ oznacza wyznacznik macierzy powstałej z podstawienia
za $i$-tą kolumnę wektora wyrazów wolnych $b$:
$$
    W_{i}
    =
    \left|
    \begin{array}{cccccc}
        a_{11}  &   a_{12}  &   \cdots  &   b_{1}   &   \cdots  &   a_{1n}\\
        a_{21}  &   a_{22}  &   \cdots  &   b_{2}   &   \cdots  &   a_{2n}\\
        \vdots  &   \vdots  &   \ddots  &   \vdots  &   \ddots  &   \vdots \\
        a_{n1}  &   a_{n2}  &   \cdots  &   b_{n}   &   \cdots  &   a_{nn}\\
    \end{array}
    \right|.
$$

%%%%%%%%%%%%%%%%%%%%%%%%%%%%%%%%%%%%%%%%%%%%%%%%%%%%%%%%%%%%%%%%
\section{Twierdzenie Kroneckera-Capelliego}
Układ równań $Ax = b$ posiada rozwiązanie wtedy i tylko wtedy,
gdy rząd macierzy $A$ jest równy rzędowi macierzy uzupełnionej
$B = [A|b]$.
Przy czym, układ ten ma:
\begin{itemize}
    \item dokładnie jedno rozwiązanie, gdy $rk(A) = rk(B) = n$
        ($n$ - liczba niewiadomych),
    \item nieskończenie wiele rozwiązań zależnych od $n-r$
        parametrów, gdy $rk(A) = rk(B) = r < n$.
\end{itemize}

%%%%%%%%%%%%%%%%%%%%%%%%%%%%%%%%%%%%%%%%%%%%%%%%%%%%%%%%%%%%%%%%
\section{Metoda Eliminacji Gaussa}
\begin{itemize}
    \item Załóżmy, że $a_{11} \neq 0$ - jeśli $a_{11} = 0$, to
        to zmieniamy numerację zmiennych.
    \item Dla $i$-tego wiersza przemnażamy pierwsze równanie przez
        $\frac{a_{i1}}{a_{11}}$ i odejmujemy od $i$-tego wiersza.
    \item Otrzymujemy wówczas:
    $$
        \left\{
        \begin{array}{ccccccccc}
            a_{11}x_{1} &   +   &   a_{12}x_{2}     &   +   &   \dots   &   +   &   a_{1n}x_{n}     &   =   &   b_{1}\\
                        &       &   a_{22}^{1}x_{2} &   +   &   \dots   &   +   &   a_{2n}^{1}x_{n} &   =   &   b_{2}^{1}\\
            \vdots      &       &   \vdots          &       &   \vdots  &       &   \vdots          &       &   \vdots\\
                        &       &   a_{m2}^{1}x_{2} &   +   &   \dots   &   +   &   a_{mn}^{1}x_{n} &   =   &   b_{m}^{1}\\
        \end{array}
        \right.
    $$
    gdzie:
    \begin{itemize}
        \item $a_{ij}^{1} = a_{ij} - \frac{a_{i1}}{a_{11}}a_{1j}$,
        \item $b_{i}^{1} = b_{i} - \frac{a_{i1}}{a_{11}}b_{1}$.
    \end{itemize}
    \item Postępujemy analogicznie dla kolejnych wierszy,
        aż otrzymamy układ w postaci:
    $$
        \left\{
        \begin{array}{cccccccccc}
            a_{11}x_{1} &   +   &   a_{12}x_{2}     &   +   &   a_{13}x_{3}     &   \dots   &   +   &   a_{1n}x_{n}         &   =   &   b_{1}\\
                        &       &   a_{22}^{1}x_{2} &   +   &   a_{23}^{1}x_{3} &   \dots   &   +   &   a_{2n}^{1}x_{n}     &   =   &   b_{2}^{1}\\
                        &       &                   &       &   a_{33}^{2}x_{3} &   \dots   &   +   &   a_{3n}^{2}x_{n}     &   =   &   b_{3}^{2}\\
                        &       &                   &       &                   &           &       &   \vdots              &       &   \vdots\\
                            &           &       &   a_{mn}^{(r-1)}x_{r}     &   +   & \cdots &  +   &   a_{rn}^{(r-1)}x_{n} &   =   &   b_{r}^{(r-1)}\\
        \end{array}
        \right.
    $$
\end{itemize}

\end{document}
